% Einheit 3 von 3 – Produktregel (bank/ableitungsregeln/e3.jsonl, zusammenbau v0.1)
%% TODO Zeitmarke Einheit 3: Marken-Zeile nicht lesbar
%% TODO Zweigzeile Teil 1: was hier gelernt wird (Satz in Schülersprache aus dem Einheitstitel) – Einheit 3
\einheitenkopf[e3]{Einheit 3 von 3 · Produktregel}
\zweigzeile{Abitur GK}
\setcounter{aufgabe}{16}

%% TODO Titel: Ich-kann-Satz fehlt in der Bank (Platzhalter: Kettenname) – Nr. 17
%% TODO Anweisung: ein Satz über der ersten Teilaufgabe fehlt in der Bank – Nr. 17
\begin{aufgabe}{Produktregel}
%% TODO \rechenplatz{n}: Schrittzahl des Grundfalls fehlt in der Bank (nur bei fester Schreibform, 2.3 b)
\begin{teile}
\teil $f(x) = x^3 \cdot e^x$ – welche Regeln? \\ \kreuz{nur Produktregel} \\ \kreuz{Produkt- und Kettenregel} \\ \kreuz{nur Kettenregel}
\teil $f(x) = x \cdot e^x$ – $f'(x)$, $e^x$ ausgeklammert? $f'(x) =$ \leerfeld
\teil Die Zahl der Neuanmeldungen bei einer App in Tausend pro Tag wird durch $f(t) = 5t \cdot e^{-\frac{t}{20}}$ beschrieben, $t$ in Tagen nach dem Start. Weise nach: $f'(t) = 5 \cdot (1 - \frac{t}{20}) \cdot e^{-\frac{t}{20}}$. \hfill (Abitur 2021 LK)
\teil $f(x) = 3 \cdot (x^2 - 4) \cdot e^x$. Weise nach: $f'(x) = 3 \cdot (x^2 + 2x - 4) \cdot e^x$. \hfill (Abitur 2024 GK)
\teil $f(x) = (x^2 - 4x + 1) \cdot e^{-x}$. Weise nach: $f'(x) = -(x - 1) \cdot (x - 5) \cdot e^{-x}$, und gib die möglichen Extremstellen ohne weitere Rechnung an. \hfill (Abitur 2024 LK) \\ $x_1 =$ \leerfeld , $x_2 =$ \leerfeld
\teil Gegeben ist $f(x) = (x^2 + 2x + 3) \cdot e^{-x}$ mit $f'(x) = (-x^2 - 1) \cdot e^{-x}$. Weise nach: $f''(x) = (x - 1)^2 \cdot e^{-x}$, und folgere daraus etwas über die Krümmung des Graphen. \hfill (Abitur 2018 LK)
\teil Die Abbildung zeigt den Graphen der linearen Funktion $g$ und den Graphen von $f$ mit dem Tiefpunkt $T$. Für $h(x) = f(x) \cdot g(x)$: Steigung der Tangente an den Graphen von $h$ an der Stelle $x = 2$? \hfill (Abitur 2026 GK) \\

\begin{ksys}[xmin=-1,xmax=6,ymin=-2,ymax=6,ablesen] \gerade{-1}{5}{g} \parabel{1}{2}{-1}{f} \tiefpunkt{2}{-1}{T} \end{ksys}
$h'(2) =$ \leerfeld
\teil Die Funktionen $f$ und $g$ sind in $\mathbb{R}$ definiert und differenzierbar mit $g(x) = f(x) \cdot e^{-x}$. Weise nach: Hat der Graph von $g$ im Punkt $(a|g(a))$ eine waagerechte Tangente, dann gilt $f'(a) = f(a)$. \hfill (Abitur 2025 LK)
\end{teile}
\end{aufgabe}

%% TODO Titel: Ich-kann-Satz fehlt in der Bank (Platzhalter: Kettenname) – Nr. 18
%% TODO Anweisung: ein Satz über der ersten Teilaufgabe fehlt in der Bank – Nr. 18
\begin{aufgabe}{Produktregel – Fehler finden, Begründen}
\begin{teile}
\teil Ich finde den Fehler: Zu $f(x) = x^4 \cdot e^x$ schreibt Ole $f'(x) = 4x^3 \cdot e^x$.
\teil Begründe, warum die Gleichung $(x - 7) \cdot e^x = 0$ nur die Lösung $x = 7$ hat.
\end{teile}
\end{aufgabe}
